% 原书分章 PDF 第 17—20 页：第 25.5 节及定理 25.4 的完整证明。
\section{数据自适应带宽选择器结果的推广}

令 $\Psi:\mathcal M\to\mathbb R$ 为感兴趣且不具路径可微性的目标参数。
令 $\Psi_h:\mathcal M\to\mathbb R$，$h\in\mathbb R_{\geq0}$，
为逼近 $\Psi$ 的路径可微参数族。这里的逼近是指，
$\{\Psi_h(P_0)\in\mathbb R:h\in\mathbb R_{\geq0}\}$
构成 $\Psi(P_0)=\lim_{h\to0}\Psi_h(P_0)$ 的一族路径可微逼近。
本章此前假定该参数族由核平滑器生成，而本节允许用一般方式定义这一目标参数族。
由下式定义 $R_{20,h}(P,P_0)$：
\[
 \Psi_h(P)-\Psi_h(P_0)=-P_0D_h^*(P)+R_{20,h}(P,P_0),
\]
其中 $D_h^*(P)$ 为 $\Psi_h:\mathcal M\to\mathbb R$ 在 $P$ 处的典范梯度。
假定对某个 $\tilde D_h(P_0)$，有
$D_h^*(P_0)=h^{-\alpha_0}\tilde D_h(P_0)$，并且当 $h\to0$ 时，
对某个 $\sigma_0^2$，有 $E_{P_0}\{\tilde D_h(P_0)\}^2\to\sigma_0^2>0$。
也就是说，对于某个 $\alpha_0>0$，典范梯度的方差的行为如同 $h^{-2\alpha_0}$。
此外，关于该参数族的偏差，令
$\Psi_h(P_0)-\Psi(P_0)=h^{J_0}B_0(J_0)+o(h^{J_0})$，
其中 $J_0>0$，$B_0(J_0)$ 为某一项。
我们允许 $\alpha_0$ 与 $J_0$ 均未知，且可依赖真实数据分布 $P_0$。
此前各节采用基于 $d$ 维核的 $\Psi$ 的核平滑逼近，
其中 $\alpha_0=2d$ 已知，而 $J_0$ 表示 $\Psi(P_0)$ 在 $x_0$ 处未知的底层光滑性。

令 $\sigma_{0h}^2\equiv E_{P_0}\{D_h^*(P_0)\}^2$
为典范梯度的方差。由上述假设，有
$\sigma_{0h}^2\approx h^{-2\alpha_0}\sigma_0^2$。
令 $\sigma_{nh}^2$ 为 $\sigma_{0h}^2$ 的相合估计量，使当 $n\to\infty$ 时，
$\sigma_{nh}^2/\sigma_{0h}^2$ 依概率趋于 1。
令 $r_h(n)\equiv(n/\sigma_{hn}^2)^{0.5}$
表示 $\Psi_h(P_0)$ 的 CV-TMLE 的标准化速率。
注意，$r_h(n)$ 的行为如同 $h^{\alpha_0}n^{0.5}$。
令 $\psi_{nh}^*$ 为 CV-TMLE，并假定在适当条件下，
对任意序列 $h_n\to0$，当 $n\to\infty$ 时，
\[
 r_{h_n}(n)(\psi^*_{nh_n}-\psi_{0h_n})\Rightarrow_dN(0,1).
\]
可用与上一节分析 CV-TMLE 类似的方法建立这一结果。
因此，$\psi_{nh}^*$ 的方差的行为如同 $h^{-2\alpha_0}n^{-1}$，
偏差的行为如同 $h^{J_0}$，从而关于均方误差（MSE）的最优带宽速率为
\[
 h_n^*=n^{-1/2(J_0+\alpha_0)}.
\]
现在，我们希望与上一节的提议及分析类似，提出数据自适应选择器 $h_n$，
并建立其行为如同 $h_n^*$ 的结果。
与上一节一样，我们将以一步估计量为基础构造选择器 $h_n$：
\[
 \psi_{nh}^1=E_{B_n}\Psi_h(\mathbf P_{n,B_n})
 +E_{B_n}P^1_{n,B_n}D_h^*(\mathbf P_{n,B_n}),
\]
其中 $\mathbf P_{n,B_n}\in\mathcal M$ 为基于训练样本
$\{O_i:B_n(i)=0\}$ 的 $P_0$ 估计量。
对于用户给定的某个序列 $C_n$，令 $h_n$ 为
\[
 \frac{d}{dh_n}\psi^1_{nh_n}
 +C_n\frac{d}{dh_n}\sigma_{nh_n}n^{-0.5}=0
\]
的一个解。关于解的存在性与 $C_n$ 的选择，参见上一节的讨论；
$C_n$ 或者关于 $n$ 一致有界，或者非常缓慢地趋于无穷。

\begin{bookresult}{定理}{25.4}
作如下假设。
\begin{itemize}
\item 对某个 $\tilde D_{1h}(P)$，
$\frac{d}{dh}D_h^*(P)=h^{-1}\tilde D_{1h}(P)$，
其中当 $h\to0$ 时，对某个 $\sigma_{10}^2>0$，有
$h^{2\alpha_0}P_0\{\tilde D_{1h}(P_0)\}^2\to\sigma_{10}^2$。
换言之，$\tilde D_{1h}(P_0)$ 作为 $h$ 的函数，其行为如同 $D_h^*(P_0)$。
\item 对任意序列 $h\to0$，有
\[
\begin{aligned}
 E_{B_n}(P^1_{n,B_n}-P_0)\tilde D_{1h_n}(\mathbf P_{n,B_n})
 &=Z_n/(n^{0.5}h_n^{\alpha_0})\\
 &\quad+o_P(n^{-0.5}h_n^{-\alpha_0}),
\end{aligned}
\]
其中 $Z_n\Rightarrow_dN(0,\sigma_{01}^2)$。
\item
\[
 h_nE_{B_n}\frac{d}{dh}R_{20,h_n}(\mathbf P_{n,B_n},P_0)
 =o_P(n^{-0.5}h_n^{-\alpha_0}).
\]
\item
\[
 -\frac{d}{dh}\sigma_{nh}
 =h^{-(\alpha_0+1)}c_{nh}+o_P(h^{-(\alpha_0+1)}),
\]
其中对于某个 $\delta>0$，$c_{nh}>0$ 关于 $h\in(0,\delta)$ 一致远离零与无穷。
\item 对任意序列 $h\to0$，有
\[
 -\frac{d}{dh}\Psi_h(P_0)
 \sim c(J_0,h)h^{J_0-1}+o(h^{J_0-1}),
\]
其中对于某个 $\delta>0$，$c(J_0,h)$ 关于 $h\in(0,\delta)$ 一致远离零与无穷。
\end{itemize}
则
\[
 c(J_0,h_n)h_n^{J_0}+o(h_n^{J_0})
 =n^{-0.5}h_n^{-\alpha_0}
 \{Z_n-c_{n,h_n}C_n+o_P(1)\}.
\]
这证明，若 $\limsup_nC_n<\infty$，则 $h_n$ 通过平衡偏差平方与方差达到最优速率，
且一般而言，
\[
 h_n\sim n^{-1/2(J_0+\alpha_0)}C_n^{1/(J_0+\alpha_0)}.
\]
\end{bookresult}

由于 $D_h^*(P_0)=h^{-\alpha_0}\tilde D_h(P_0)$，其中
$P_0\{\tilde D_h(P_0)\}^2\to\sigma_0^2$，
关于 $\frac{d}{dh}\sigma_{nh}$ 的第三条假设是合理的。
类似地，由于
$\Psi_h(P_0)-\Psi(P_0)=h^{J_0}B_0(J_0)+o(h^{J_0})$，
第四条假设也是合理的。

\begin{proof}
我们将按照与定理 25.3 的证明类似的方法分析数据自适应选择器 $h_n$。
有
\[
\begin{aligned}
 \psi^1_{nh_n}-\Psi_{h_n}(P_0)
 &=E_{B_n}(P^1_{n,B_n}-P_0)D_{h_n}^*(\mathbf P_{n,B_n})\\
 &\quad+E_{B_n}R_{20,h_n}(\mathbf P_{n,B_n},P_0).
\end{aligned}
\]
关于 $h$ 在 $h=h_n$ 处求导，并利用 $h_n$ 的定义，得到
\[
\begin{aligned}
 C_n\frac{d}{dh_n}\sigma_{nh_n}/n^{0.5}
 -\frac{d}{dh_n}\Psi_{h_n}(P_0)
 &=E_{B_n}(P^1_{n,B_n}-P_0)
 \frac{d}{dh_n}D_{h_n}^*(\mathbf P_{n,B_n})\\
 &\quad+E_{B_n}\frac{d}{dh_n}
 R_{20,h_n}(\mathbf P_{n,B_n},P_0).
\end{aligned}
\]
我们假定 $\frac{d}{dh}D_h^*(P)=h^{-1}\tilde D_{1h}(P)$，
其中当 $h\to0$ 时，
$h^{2\alpha_0}P\{\tilde D_{1h}(P)\}^2\to\sigma_1^2$。
我们假定在正则条件下，有下式；这些条件类似于证明一步估计量
$\psi^1_{nh}$ 渐近正态性所需的条件：
\[
 E_{B_n}(P^1_{n,B_n}-P_0)\tilde D_{1h_n}(\mathbf P_{n,B_n})
 =Z_n/(n^{0.5}h_n^{\alpha_0})+o_P(1),
\]
其中 $Z_n\Rightarrow_dN(0,\sigma_{01}^2)$。
此外，我们假定
\[
 h_nE_{B_n}\frac{d}{dh_n}R_{20,h_n}(\mathbf P_{n,B_n},P_0)
 =o_P(n^{-0.5}h_n^{-\alpha_0}).
\]
由此得到
\[
\begin{aligned}
 -\frac{d}{dh_n}\Psi_{h_n}(P_0)
 &=h_n^{-1}Z_nn^{-0.5}h_n^{-\alpha_0}
 +h_n^{-1}o_P(n^{-0.5}h_n^{-\alpha_0})\\
 &\quad-C_nn^{-0.5}\frac{d}{dh_n}\sigma_{nh_n}.
\end{aligned}
\]
我们还假定
\[
 -\frac{d}{dh}\sigma_{nh}=h^{-(\alpha_0+1)}c_{nh},
\]
其中序列 $c_{nh}>0$ 远离 0 与 $\infty$。于是得到
\[
\begin{aligned}
 -\frac{d}{dh_n}\Psi_{h_n}(P_0)
 &=Z_nn^{-0.5}h_n^{-(\alpha_0+1)}
 +o_P(n^{-0.5}h_n^{-(\alpha_0+1)})\\
 &\quad-C_nn^{-0.5}c_{nh_n}h_n^{-(\alpha_0+1)}.
\end{aligned}
\]
我们假定
\[
 -\frac{d}{dh}\Psi_h(P_0)
 \sim c(J_0,h)h^{J_0-1}+o(h^{J_0-1}),
\]
其中对某个 $\delta>0$，$c(J_0,h)$ 关于 $h\in(0,\delta)$
一致远离零与无穷。将其代入最后一个等式，得到
\[
 c(J_0,h_n)h_n^{J_0}+o(h_n^{J_0})
 =n^{-0.5}h_n^{-\alpha_0}
 \{Z_n-c_{n,h_n}C_n+o_P(1)\}.
\]
这证明，若 $\limsup_nC_n<\infty$，则 $h_n$ 平衡偏差平方与方差，
且一般而言，
\[
 h_n\sim n^{-1/2(J_0+\alpha_0)}C_n^{1/(J_0+\alpha_0)}.
\]
\end{proof}
